% ---------------------------------------------------------------------
%  CHƯƠNG 1 — MỞ ĐẦU
% ---------------------------------------------------------------------
\chapter{Mở đầu}
\label{ch:modau}

\begin{muctieu}
Chương mở đầu trình bày lý do lựa chọn đề tài, mục tiêu cần đạt, đối tượng,
phạm vi thực hiện và bố cục của báo cáo. Qua đó, người đọc có thể nắm được
bức tranh tổng quan về ứng dụng trí tuệ nhân tạo trong nhận diện quần áo,
gợi ý phối đồ và thử trang phục ảo.
\end{muctieu}

\section{Lý do chọn đề tài}

Trong đời sống hiện nay, nhu cầu lựa chọn trang phục phù hợp với hoàn cảnh,
màu sắc và phong cách cá nhân ngày càng được quan tâm. Tuy nhiên, nhiều người
gặp khó khăn trong việc xác định loại quần áo, lựa chọn những món đồ phù hợp
với nhau hoặc hình dung trang phục sẽ trông như thế nào khi mặc. Việc thử
trực tiếp nhiều bộ trang phục cũng tốn thời gian, đặc biệt khi mua sắm trực
tuyến hoặc quản lý số lượng lớn quần áo cá nhân.

Sự phát triển của trí tuệ nhân tạo, đặc biệt là thị giác máy tính
\en{(computer vision)} và học sâu \en{(deep learning)}, cho phép máy tính
phân tích hình ảnh và nhận diện các đối tượng trong ảnh với độ chính xác
ngày càng cao. Những công nghệ này có thể được ứng dụng để nhận diện loại
quần áo, trích xuất một số đặc điểm như màu sắc và kiểu dáng, từ đó hỗ trợ
người dùng lựa chọn trang phục phù hợp.

Bên cạnh đó, điện thoại thông minh đã trở thành thiết bị phổ biến và thuận
tiện để chụp ảnh, lưu trữ thông tin cũng như sử dụng các dịch vụ thông minh.
Vì vậy, việc xây dựng một ứng dụng di động tích hợp trí tuệ nhân tạo có thể
giúp người dùng nhận diện quần áo, quản lý tủ đồ, nhận gợi ý phối đồ và thử
trang phục ảo ngay trên điện thoại.

Từ những nhu cầu thực tế trên, đề tài ``Xây dựng hệ thống gợi ý phối đồ''
được lựa chọn. Trong phạm vi đồ án, hệ thống được định hướng dưới dạng ứng
dụng di động Flutter, kết hợp nhận diện quần áo, quản lý tủ đồ, gợi ý phối
đồ và thử trang phục ảo ở mức cơ bản. Đề tài không chỉ có ý nghĩa thực tiễn
mà còn tạo điều kiện vận dụng các kiến thức về trí tuệ nhân tạo, xử lý ảnh,
phát triển ứng dụng di động, xây dựng API và quản lý dữ liệu.

\section{Mục tiêu của đồ án}

Mục tiêu tổng quát của đồ án là xây dựng một ứng dụng di động có khả năng
nhận diện quần áo từ hình ảnh, đề xuất cách phối đồ và hỗ trợ người dùng thử
trang phục ảo ở mức cơ bản.

Các mục tiêu cụ thể của đồ án gồm:

\begin{itemize}[leftmargin=1.6em, itemsep=3pt]
  \item Xây dựng mô hình trí tuệ nhân tạo có khả năng phân loại một số nhóm
        quần áo phổ biến từ ảnh do người dùng chụp hoặc tải lên.

  \item Đánh giá và so sánh ít nhất hai mô hình học sâu phù hợp với bài toán
        nhận diện hình ảnh, từ đó lựa chọn mô hình có sự cân bằng giữa độ
        chính xác và tốc độ xử lý.

  \item Hướng tới độ chính xác nhận diện từ 80\% trở lên trên tập kiểm thử
        và thời gian phản hồi không quá 3 giây trong điều kiện kiểm thử ổn
        định.

  \item Xây dựng chức năng gợi ý tối đa ba cách phối đồ phù hợp dựa trên loại
        trang phục, màu sắc và phong cách, đồng thời hạn chế các kết quả bị
        trùng lặp.

  \item Xây dựng chức năng thử trang phục ảo ở mức cơ bản bằng cách đặt ảnh
        quần áo có nền trong suốt lên ảnh người dùng và cho phép điều chỉnh
        vị trí, kích thước và góc xoay.

  \item Phát triển ứng dụng di động bằng Flutter, kết nối với hệ thống xử lý
        trí tuệ nhân tạo thông qua API.

  \item Xây dựng các chức năng hỗ trợ như đăng ký, đăng nhập, quản lý tủ đồ
        cá nhân, lưu trang phục yêu thích và lưu kết quả phối đồ.
\end{itemize}

\section{Phạm vi và đối tượng}

\subsection{Đối tượng nghiên cứu}

Đối tượng nghiên cứu của đồ án gồm:

\begin{itemize}[leftmargin=1.6em, itemsep=3pt]
  \item Các hình ảnh chứa một trang phục chính như áo thun, áo sơ mi, áo
        khoác, quần dài, quần ngắn, váy hoặc đầm.

  \item Các mô hình học sâu có khả năng phân loại hình ảnh quần áo và phù
        hợp với hệ thống có tài nguyên tính toán giới hạn.

  \item Các nguyên tắc phối hợp trang phục dựa trên loại quần áo, màu sắc
        và phong cách cơ bản.

  \item Phương pháp xác định vị trí cơ thể để đặt lớp ảnh trang phục lên ảnh
        người dùng.

  \item Quy trình phát triển ứng dụng di động bằng Flutter và kết nối ứng
        dụng với dịch vụ xử lý trí tuệ nhân tạo.
\end{itemize}

\subsection{Phạm vi thực hiện}

Trong phạm vi đồ án, hệ thống tập trung vào các chức năng sau:

\begin{itemize}[leftmargin=1.6em, itemsep=3pt]
  \item Nhận ảnh quần áo từ camera hoặc thư viện ảnh của điện thoại.

  \item Nhận diện một số loại quần áo phổ biến đã được xác định trong bộ dữ
        liệu huấn luyện.

  \item Trả về tên loại quần áo dự đoán và độ tin cậy của kết quả.

  \item Gợi ý ba phương án phối đồ dựa trên các thuộc tính cơ bản của trang
        phục và dữ liệu trong tủ đồ cá nhân.

  \item Hỗ trợ thử trang phục ảo bằng kỹ thuật ghép ảnh hai chiều
        \en{(2D overlay)}, không sử dụng mô hình sinh ảnh.

  \item Cho phép người dùng điều chỉnh thủ công vị trí, kích thước và góc
        xoay của trang phục khi thử đồ.

  \item Lưu thông tin người dùng, hình ảnh quần áo, tủ đồ và các kết quả yêu
        thích bằng Firebase.
\end{itemize}

Đồ án chưa tập trung vào việc nhận diện đồng thời nhiều loại trang phục
trong một ảnh phức tạp, xác định thương hiệu, chất liệu hoặc kích cỡ chính
xác của quần áo. Chức năng thử trang phục ảo chỉ mô phỏng bằng lớp ảnh hai
chiều nên chưa thể tái tạo chính xác độ rủ, nếp gấp, ánh sáng hoặc sự thay đổi
hình dạng của trang phục theo cơ thể người dùng.

Ứng dụng được xây dựng chủ yếu để thử nghiệm trên thiết bị Android. Hệ thống
chưa hướng tới triển khai thương mại, chưa tích hợp chức năng thanh toán và
chưa thay thế hoàn toàn việc thử trang phục trực tiếp.

\section{Bố cục báo cáo}

Phần còn lại của báo cáo được tổ chức như sau. Chương~\ref{ch:coso} trình
bày cơ sở lý thuyết về trí tuệ nhân tạo, thị giác máy tính, mô hình nhận diện
ảnh và các công nghệ được sử dụng. Chương~\ref{ch:thietke} phân tích yêu cầu
và thiết kế kiến trúc, dữ liệu, chức năng của hệ thống. Chương~\ref{ch:caidat}
mô tả quá trình xây dựng mô hình AI, phát triển ứng dụng Flutter và tích hợp
các thành phần. Chương~\ref{ch:ketqua} trình bày kết quả thực nghiệm, kiểm
thử và đánh giá hệ thống. Cuối cùng, Chương~\ref{ch:ketluan} tổng kết những
kết quả đã đạt được, nêu các hạn chế và đề xuất hướng phát triển trong tương
lai.
